% Sections 6.1 to 6.5, from lectures/L06/appendix/a_theory.md.

\section{Definition}
\label{c6:sec:definition}
En \textbf{andragradsekvation} (kvadratisk ekvation) är en ekvation på formen
\[
  ax^2 + bx + c = 0, \quad a \neq 0.
\]
Koefficienterna $a$, $b$ och $c$ är reella tal. Ekvationen kan ha \textbf{0, 1 eller 2 reella
rötter}.

\begin{exempel}
$2x^2 - 3x - 5 = 0$ är en andragradsekvation med $a = 2$, $b = -3$ och $c = -5$.
\end{exempel}

\section{Diskriminant}
\label{c6:sec:diskriminant}
\textbf{Diskriminanten} $D$ avgör antalet reella rötter:
\[
  D = b^2 - 4ac
\]

\begin{narrowtable}{@{}ll@{}}
  \thead{Villkor} & \thead{Antal rötter} \\
  \midrule
  $D > 0$ & Två skilda reella rötter \\
  $D = 0$ & En dubbelrot (en rot med multipliciteten 2) \\
  $D < 0$ & Inga reella rötter \\
\end{narrowtable}

\section{Lösningsmetoder}
\label{c6:sec:losningsmetoder}

\subsection{Faktoriseringsmetoden}
\label{c6:sec:faktorisering}
Om $ax^2 + bx + c$ kan faktoriseras till $a(x - r_1)(x - r_2)$ är rötterna $x = r_1$ och
$x = r_2$.

\begin{exempel}
Lös $x^2 - 5x + 6 = 0$. Faktorisera:
\[
  (x - 2)(x - 3) = 0
\]
Rötterna ges av $x - 2 = 0$ eller $x - 3 = 0$:
\[
  x_1 = 2, \quad x_2 = 3
\]
\end{exempel}

\begin{exempel}[Konjugatregeln]
Lös $x^2 - 9 = 0$.
\[
  (x + 3)(x - 3) = 0 \quad \Rightarrow \quad x = -3 \text{ eller } x = 3
\]
\end{exempel}

\subsection{PQ-formeln}
\label{c6:sec:pq}
PQ-formeln gäller för ekvationer på \emph{normerad} form $x^2 + px + q = 0$, det vill säga med
$a = 1$:
\[
  x = -\frac{p}{2} \pm \sqrt{\left(\frac{p}{2}\right)^2 - q}
\]

\begin{exempel}
Lös $x^2 - 4x + 3 = 0$. Här är $p = -4$ och $q = 3$:
\begin{gather*}
  x = \frac{4}{2} \pm \sqrt{\left(\frac{4}{2}\right)^2 - 3} = 2 \pm \sqrt{4 - 3} = 2 \pm 1 \\*
  x_1 = 3, \quad x_2 = 1
\end{gather*}
\end{exempel}

\subsection{ABC-formeln (kvadratiska formeln)}
\label{c6:sec:abc}
ABC-formeln gäller generellt för $ax^2 + bx + c = 0$:
\[
  x = \frac{-b \pm \sqrt{b^2 - 4ac}}{2a}
\]

\begin{exempel}
Lös $2x^2 + 5x - 3 = 0$. Här är $a = 2$, $b = 5$ och $c = -3$:
\[
  D = 5^2 - 4 \cdot 2 \cdot (-3) = 25 + 24 = 49
\]
\[
  x = \frac{-5 \pm \sqrt{49}}{4} = \frac{-5 \pm 7}{4}
  \quad \Rightarrow \quad x_1 = 0{,}5, \quad x_2 = -3
\]
\end{exempel}

\subsection{Kvadratkomplettering}
\label{c6:sec:kvadratkomplettering}
Varje andragradsekvation kan skrivas på formen $(x + k)^2 = m$ genom att ”komplettera till en
kvadrat”. För $x^2 + bx + c = 0$ går det till så här:
\begin{enumerate}
  \item Flytta konstanten: $x^2 + bx = -c$.
  \item Lägg till $\left(\frac{b}{2}\right)^2$ på båda sidor.
  \item Vänsterledet blir en perfekt kvadrat.
\end{enumerate}

\begin{exempel}
Lös $x^2 + 6x + 5 = 0$.
\[
  x^2 + 6x = -5
\]
Lägg till $\left(\frac{6}{2}\right)^2 = 9$:
\[
  x^2 + 6x + 9 = -5 + 9 = 4
\]
\[
  (x + 3)^2 = 4 \quad \Rightarrow \quad x + 3 = \pm 2
  \quad \Rightarrow \quad x_1 = -1, \quad x_2 = -5
\]
\end{exempel}

\section{Vietas formler}
\label{c6:sec:vieta}
Om $x_1$ och $x_2$ är rötterna till $x^2 + px + q = 0$ gäller
\begin{align*}
  x_1 + x_2 &= -p \qquad \text{(summan av rötterna)}, \\
  x_1 \cdot x_2 &= q \qquad \text{(produkten av rötterna)}.
\end{align*}

\begin{exempel}[Snabbkontroll]
I \secref{c6:sec:pq} gav $x^2 - 4x + 3 = 0$ rötterna $x_1 = 3$ och $x_2 = 1$. Vietas formler
bekräftar dem:
\[
  3 + 1 = 4 = -(-4)\ \checkmark, \qquad 3 \cdot 1 = 3\ \checkmark
\]
\end{exempel}

\section{Tillämpning i elektroteknik}
\label{c6:sec:tillampning}

\subsection{Ström från effekt och resistans}
\label{c6:sec:strom}
Effekten i ett motstånd $R$ ges av $P = RI^2$. Givet $P$ och $R$ kan strömmen $I$ beräknas:
\[
  RI^2 = P \quad \Rightarrow \quad I^2 = \frac{P}{R} \quad \Rightarrow \quad I = \sqrt{\frac{P}{R}}
\]
Det är en \emph{ren} andragradsekvation, utan linjär term.

\begin{exempel}
Med $P = 8\,\text{W}$ och $R = 50\,\Omega$ blir
\[
  I = \sqrt{\frac{8}{50}} = \sqrt{0{,}16} = 0{,}4\,\text{A}.
\]
\end{exempel}

\subsection{Resistans från seriekoppling}
\label{c6:sec:serie}
Två motstånd $R_1$ och $R_2$ är seriekopplade med $R_1 + R_2 = 10\,\Omega$ och
$R_1 \cdot R_2 = 24\,\Omega^2$. Enligt Vietas formler är de rötterna till
\[
  x^2 - 10x + 24 = 0 \quad \Rightarrow \quad x = \frac{10 \pm 2}{2},
\]
alltså
\[
  R_1 = 6\,\Omega, \quad R_2 = 4\,\Omega.
\]
